% Bagian: turing-machines
% Bab: undecidability
% Subbagian: halting-problem

\documentclass[../../../include/open-logic-section]{subfiles}

\begin{document}

\olfileid[id]{tur}{und}{hal}
\olsection{Masalah Penghentian}

\begin{explain}
Andaikan kita telah menetapkan suatu enumerasi deskripsi mesin Turing. Dengan
demikian, setiap mesin Turing memperoleh suatu \emph{indeks}: tempatnya dalam
enumerasi $M_1$, $M_2$, $M_3$, \dots{} dari deskripsi mesin Turing.

Kita mengetahui bahwa pasti ada fungsi yang tidak dapat dihitung oleh mesin
Turing: himpunan deskripsi mesin Turing---dan karenanya himpunan mesin
Turing---bersifat !!{enumerable}, tetapi himpunan semua fungsi dari $\Nat$ ke
$\Nat$ tidak demikian. Namun, kita juga dapat menemukan contoh tertentu dari
fungsi yang tidak dapat dihitung. Salah satunya ialah fungsi penghentian.
\end{explain}

\begin{defn}[Fungsi penghentian]
 \emph{Fungsi penghentian}~$h$ didefinisikan sebagai
\[
h(e,n) =
\begin{cases}
  \text{0} & \text{jika mesin~$M_e$ tidak berhenti pada masukan $n$} \\
  \text{1} & \text{jika mesin~$M_e$ berhenti pada masukan $n$}
\end{cases}
\]
\end{defn}

\begin{defn}[Masalah penghentian]
\emph{Masalah Penghentian} adalah masalah menentukan (untuk setiap $e$, $n$)
apakah mesin Turing~$M_e$ berhenti pada masukan berupa $n$ simbol~$\TMstroke$.
\end{defn}

\begin{explain}
Kita menunjukkan bahwa~$h$ tidak dapat dihitung oleh mesin Turing dengan
menunjukkan bahwa fungsi terkait,~$s$, tidak dapat dihitung oleh mesin Turing.
Bukti ini bergantung pada kenyataan bahwa apa pun yang dapat dihitung oleh
mesin Turing dapat dihitung oleh mesin Turing disiplin
(\olref[mac][dis]{sec}), dan kenyataan bahwa dua mesin Turing dapat dirangkai
untuk membuat satu mesin (\olref[mac][cmb]{sec}).
\end{explain}

\begin{defn} Fungsi~$s$ didefinisikan sebagai
\[
s(e) =
\begin{cases}
  \text{0} & \text{jika mesin~$M_e$ tidak berhenti pada masukan $e$} \\
  \text{1} & \text{jika mesin~$M_e$ berhenti pada masukan $e$}
\end{cases}
\]
\end{defn}

\begin{lem}
Fungsi~$s$ tidak dapat dihitung oleh mesin Turing.
\end{lem}

\begin{proof}
Andaikan, demi memperoleh kontradiksi, bahwa fungsi~$s$ dapat dihitung oleh
mesin Turing. Dengan demikian, akan ada mesin Turing~$S$ yang menghitung~$s$.
Tanpa mengurangi keumuman, kita dapat mengandaikan bahwa ketika $S$ berhenti,
mesin itu berhenti sambil memindai petak pertama (yakni, bahwa mesin itu
terdisiplin). Mesin ini dapat ``dirangkai'' dengan mesin lain~$J$, yang
berhenti jika dijalankan pada masukan~$0$ (yakni, jika mesin membaca
$\TMblank$ dalam keadaan awal ketika memindai petak di sebelah kanan simbol
ujung pita), dan selain itu bergerak terus ke kanan tanpa pernah berhenti.
$S \concat J$, yaitu mesin yang dibuat dengan merangkai $S$ dengan~$J$,
merupakan mesin Turing; jadi, mesin itu adalah $M_e$ untuk suatu~$e$ (yakni,
mesin tersebut muncul di suatu tempat dalam enumerasi). Jalankan $M_e$ pada
masukan berupa $e$ simbol~$\TMstroke$. Ada dua kemungkinan: $M_e$ berhenti
atau tidak berhenti.
\begin{enumerate}
\item Andaikan $M_e$ berhenti pada masukan berupa $e$ simbol~$\TMstroke$.
  Maka $s(e) = 1$. Jadi, ketika dijalankan pada~$e$, $S$ berhenti dengan satu
  simbol~$\TMstroke$ sebagai keluaran pada pita. Kemudian $J$ mulai bekerja
  dengan satu simbol~$\TMstroke$ pada pita. Dalam hal itu, $J$ tidak berhenti.
  Namun, $M_e$ adalah mesin $S \concat J$, sehingga seharusnya mesin itu
  melakukan persis apa yang dilakukan $S$ dan kemudian $J$ (yakni, dalam hal
  ini, bergerak terus ke kanan dan tidak pernah berhenti). Jadi, $M_e$ tidak
  mungkin berhenti pada masukan berupa $e$ simbol~$\TMstroke$.

\item Sekarang andaikan $M_e$ tidak berhenti pada masukan berupa $e$
  simbol~$\TMstroke$. Maka $s(e) = 0$, dan ketika dijalankan pada masukan~$e$,
  $S$ berhenti dengan pita kosong. Ketika dijalankan pada pita kosong, $J$
  langsung berhenti. Sekali lagi, $M_e$ melakukan apa yang dilakukan $S$ dan
  kemudian~$J$, sehingga $M_e$ harus berhenti pada masukan berupa $e$
  simbol~$\TMstroke$.
\end{enumerate}
Dalam kedua kasus, kita memperoleh kontradiksi dengan pengandaian kita. Ini
menunjukkan bahwa mesin Turing~$S$ semacam itu tidak mungkin ada: $s$ tidak
dapat dihitung oleh mesin Turing.
\end{proof}

\begin{thm}[Ketakterputusan Masalah Penghentian]
\ollabel{thm:halting-problem} Masalah penghentian tidak dapat diselesaikan;
yakni, fungsi~$h$ tidak dapat dihitung oleh mesin Turing.
\end{thm}

\begin{proof}
Andaikan $h$ dapat dihitung oleh mesin Turing, misalnya oleh mesin Turing~$H$.
Kita dapat menggunakan $H$ untuk membangun suatu mesin Turing yang
menghitung~$s$: pertama, buat salinan masukan (yang dipisahkan oleh satu
simbol~$\TMblank$). Kemudian bergerak kembali ke awal dan jalankan~$H$. Jelas
kita dapat membuat mesin yang melakukan langkah pertama tersebut (lihat
\cref{tur:mac:dis:prob:copier}), dan seandainya $H$ ada, kita dapat
``merangkainya'' dengan mesin penyalin semacam itu untuk memperoleh mesin baru
yang menentukan apakah $M_e$ berhenti pada masukan~$e$, yakni menghitung~$s$.
Namun, kita telah menunjukkan bahwa mesin semacam itu tidak mungkin ada.
Karena itu, $h$ juga tidak dapat dihitung oleh mesin Turing.
\end{proof}

\begin{prob}
Masalah Tiga-Penghentian (3-Halt) adalah masalah memberikan suatu prosedur
keputusan untuk menentukan apakah mesin Turing yang dipilih secara sembarang
berhenti pada masukan berupa tiga simbol~$\TMstroke$ pada pita yang selain itu
kosong. Buktikan bahwa masalah 3-Halt tidak dapat diselesaikan.
\end{prob}

\begin{prob}
Tunjukkan bahwa jika masalah penghentian dapat diselesaikan bagi pasangan mesin
Turing dan masukan $M_e$ dan~$n$ dengan $e \neq n$, maka masalah itu juga dapat
diselesaikan bagi kasus $e = n$.
\end{prob}

\begin{prob}
Kita telah membuktikan bahwa masalah penghentian tidak dapat diselesaikan jika
masukannya berupa bilangan~$e$, yang mengidentifikasi mesin Turing~$M_e$
melalui suatu enumerasi semua mesin Turing. Bagaimana jika kita secara langsung
mengizinkan deskripsi mesin Turing dari \olref[tur][und][enu]{sec} sebagai
masukan? Dapatkah ada mesin Turing yang memutuskan masalah penghentian tetapi
menerima deskripsi mesin Turing sebagai masukan, bukan indeks? Jelaskan mengapa
dapat atau mengapa tidak.
\end{prob}

\begin{prob} Tunjukkan bahwa fungsi \emph{parsial}~$s'$ yang didefinisikan oleh
  \[
  s'(e) =
  \begin{cases}
    \text{1} & \text{jika mesin~$M_e$ berhenti pada masukan $e$}\\
    \text{tidak terdefinisi} & \text{jika mesin~$M_e$ tidak berhenti pada masukan $e$}
  \end{cases}
  \]
  \emph{dapat} dihitung oleh mesin Turing.
\end{prob}

\end{document}
